\documentclass[10pt,a4paper]{amsart}
\usepackage[margin=19mm]{geometry}
\usepackage{amsmath,amssymb,mathtools}
\usepackage[scr=rsfs,cal=euler]{mathalfa}
\usepackage{xfrac}
\usepackage[unicode,hidelinks,bookmarks=false]{hyperref}
\newcommand{\ie}{i.e.\ }
\newcommand{\cf}{cf.\ }
\newcommand{\resp}{respectively\ }
\newcommand{\Subsection}{\subsection}
\providecommand{\nobreakdash}{\nobreak\hbox{-}}
\setlength{\emergencystretch}{2em}
\multlinegap0pt
\usepackage{bm}
\usepackage{bbm}
\usepackage{dsfont}
\usepackage{xspace}
\usepackage{cancel}
\usepackage{mathtools}
\usepackage{amsthm}
\makeatletter
\@ifundefined{thm}{\newtheorem{thm}{Thm}}{}
\@ifundefined{claim}{\newtheorem{claim}{Claim}}{}
\@ifundefined{lem}{\newtheorem{lem}[thm]{Lemma}}{}
\makeatother
\title[The high order spectral radius of graphs without long cycles]{The high order spectral radius of graphs without long cycles or paths}
\author{Yuntian Wang, Lizhu Sun, Changjiang Bu}
\date{Source: arXiv:2510.04461v1 (CC BY 4.0)}
\begin{document}
\maketitle

\noindent\emph{Source and licence.} The high order spectral radius of graphs without long cycles or paths. Version arXiv:2510.04461v1 (CC BY 4.0), distributed under the Creative Commons Attribution 4.0 International licence, as recorded in the arXiv licence metadata for this exact version and shown on the abstract page https://arxiv.org/abs/2510.04461v1. Full rights notice: licenses/arxiv-2510.04461-CC-BY-4.0.txt.\par\smallskip

\noindent\textit{Editorial scope.} Counted unit: the Lem whose source label is \texttt{CL4} in the article (source0.tex lines 398--400, source label \texttt{CL4}), delivered together with the article's own complete original proof of it (source0.tex lines 402--529, one \texttt{proof} environment of 128 lines). No mathematics is written, repaired or re-derived: statement and proof are the article's own text line for line. Required context retained uncounted: CL1 (lines 294--296). Every definition, display and label of those lines is retained unchanged. Omitted from this package and not redistributed: 1--293, 297--397, 401--401, 530--709. Nothing inside a retained excerpt is omitted or altered: every retained body is delivered exactly as the article's lines are written, so no omission receipt of any kind accompanies this package. Citations: the retained excerpts carry no citation command at all, so no bibliography entry of the article is transcribed and no reference is redistributed. Reference adaptation: none was needed and none was made; every reference inside the delivered unit resolves inside it, so no unresolved reference or citation is produced. Printed slips kept literal: no printed word, symbol, hypothesis or number of the counted statement or of its proof was altered. Layout and class: the article's own class is \texttt{elsarticle} with packages including amssymb, amsthm, amsmath, epic; the wrapper is the lane's pinned \texttt{amsart} editorial wrapper at 10pt on A4, which restates the theorem-like environments the retained text needs and adds the article's own macro definitions that the retained text needs (none), with extra wrapper packages (bm, bbm, dsfont, xspace, cancel, mathtools). The delivered numbering is the unit's own. Assets: no retained span contains a figure environment, a graphics-inclusion command, a PDF-page inclusion, a table or a TikZ picture; no image file is copied into this package, the package declares no asset dependency and no image file is redistributed with it. Licence: version arXiv:2510.04461v1 of this article is distributed under the Creative Commons Attribution 4.0 International licence, as recorded in the arXiv licence metadata for this exact version and shown on the abstract page https://arxiv.org/abs/2510.04461v1. A full rights notice is delivered at licenses/arxiv-2510.04461-CC-BY-4.0.txt. Characters: every retained excerpt is pure ASCII, so the delivered engine, XeLaTeX with its default Latin Modern text font, reports no missing character.
        \begin{lem}\label{CL1}
                Let $G$ be a graph with $n$ vertices, and let $u,v \in V(G)$, then $\rho_t(G_{u \to v}) \geq \rho_t(G)$.
        \end{lem}

        \begin{lem}\label{CL4}
                 Let $\bar{G} \in \mathrm{SPEX}_t'(n, C_{\geq k})$, $G $ is a $t$-clique connected compoment of $\bar{G} $ that satisfies $\rho_t(G) = \rho_t(\bar{G})$, then the graph $G$ does not contain any level-3 components.
        \end{lem}

        \begin{proof}

                Suppose $G$ contains a level-3 component $H$, the level-2 components among the subordinate components of $H$ be ordered by the size of their universal vertex sets in descending order as $D_1, \cdots, D_p$. We now consider the level-2 components on $H$.
            \begin{claim}\label{CC1}
                For each $h \in [p]$, the level-2 component $D_h$ has at most one subordinate component that contains a clique with more than one vertex.
            \end{claim}
            \begin{proof}
                Let $un(D_h) = m_0$. By definition, all subordinate components of the level-2 component $D_h$ are complete graphs, denote them in descending order of size as $M_1, M_2, \cdots, M_l$, and let their respective numbers of vertices be $m_1, m_2, \cdots, m_l$.
                By symmetry, for the $t$-clique Perron vector $x$ of $G$, any two vertices $i, j$ such that both belong to the same clique $M_k$, or both belong to $Un(D_h)$, have the same $t$-clique Perron vector component, hence, we denote the this $t$-clique Perron vector entry corresponding to any vertex in $M_k$ by $x_k$ and corresponding to any vertex in $Un(D_h)$ by $x_0$.
                For given $i, j \geq 1$, if $m_i > m_j$, then $x_i > x_j$, otherwise, if $m_i > m_j$ and $x_i \leq x_j$, we arbitrarily select $m_j$ vertices from $V(M_i)$ to form a induced subgraph of $M_i$, denoted as $M_i'$, then, by swapping the components of $x$ corresponding to $V(M_i')$ and $V(M_j)$, we obtain the vector $y$ with its $t$-norm equal to $1$. According to the definition, $L_r( G , M_i, x)=L_r( G ,M_j, x)$ holds for any $r$, thus, we have
                    \begin{align*}
                        &A_t(G)y^t-A_t(G)x^t
                        \\= &t\sum_{r=1}^t L_{t-r}(G, M_i, x) \left( \sum_{s=1}^r \begin{pmatrix} m_j\\s \end{pmatrix} \begin{pmatrix}m_i-m_j\\ r-s \end{pmatrix} x_i^{r-s}x_j^{s} + \begin{pmatrix}m_j\\r\end{pmatrix} x_i^r \right.
                        \\& \left.- \sum_{s=1}^r \begin{pmatrix} m_j\\s \end{pmatrix} \begin{pmatrix}m_i-m_j\\ r-s \end{pmatrix} x_i^{r} - \begin{pmatrix}m_j\\r\end{pmatrix} x_j^r \right)
                        \\= &t\sum_{r=1}^t L_{t-r}(G, M_i, x) \sum_{s=1}^{r-1} \begin{pmatrix} m_j\\s \end{pmatrix} \begin{pmatrix}m_i-m_j\\ r-s \end{pmatrix} x_i^{r-s}(x_j^{r}-x_i^{r})
                        \\\geq &0,
                    \end{align*}
                equality holds if and only if $x_i=x_j$, but since $m_i > m_j$, by comparing the components of the vertices in $M_i $ and $M_j $ through Equation (1), it can be concluded this cannot be true, therefore, we have $A_t(G)y^t-A_t(G)x^t > 0$, this contradicts the fact that $x$ is the $t$-clique Perron vector.

                If $m_2 = 1$, then the Claim already holds, so it suffices to consider the case where $m_2 \geq 2$. Suppose $m_2 \geq 2$, consider the following operation on $D_h$ to obtain a new graph $D_h'$:

                \noindent (1) Arbitrarily choose $m_2 - 1$ vertices from $V(M_1)$ to form an induced subgraph of $M_1$, denoted by $M_1'$. Connect each pair $(i, j)$, where $i \in V(M_1')$ and $j \in \bigcup_{i=2}^l V(M_i)$;

                \noindent (2) For each $i \in {2, \cdots, l}$, delete all edges in $E(M_i)$.

                It is easy to see that $D_h'$ is still a level-2 component, and it has at most one subordinate component whose clique size is greater than 1. Define a new graph $G'$ with vertex set $V(G') = V(G)$ and edge set $E(G') = E(G) \setminus E(D_h) \cup E(D_h')$, then $p(D_h') = 2m_0+m_1+m_2-1 \leq \sum_{i=0}^{m_0+1}m_i = p(D_h)$, furthermore, $c(G') \leq c(G) \leq c(\bar{G})$, and
                    \begin{align*}
                        &\rho_t(G')-\rho_t(G)
                        \\\geq & A_t(G')x^t-A_t(G)x^t
                        \\= &\sum_{r=2}^t L_{t-r}(G, D_h, x) \left( \sum_{i=2}^l m_i \sum_{s=1}^{r-1} \begin{pmatrix}m_2-1\\s\end{pmatrix} \begin{pmatrix}m_0\\r-s-1\end{pmatrix}x_1^sx_0^{r-s-1}x_i \right)\\
                        &- \sum_{r=2}^t L_{t-r}(G, D_h, x) \left( \sum_{i=2}^l \sum_{s=2}^{r} \begin{pmatrix}m_i\\s\end{pmatrix} \begin{pmatrix}m_0\\r-s\end{pmatrix}x_i^{s}x_0^{r-s} \right)
                        \\= &\sum_{r=2}^t L_{t-r}(G, D_h, x) \\
                        &\sum_{i=2}^l \sum_{s=1}^{r-1}\begin{pmatrix}m_0\\r-s-1\end{pmatrix}x_0^{r-s-1}x_i\left( \begin{pmatrix}m_2-1\\s\end{pmatrix}m_ix_1^s-\begin{pmatrix}m_i\\s+1\end{pmatrix}x_i^s \right)
                        \\\geq &0,
                    \end{align*}

                The equality in the first inequality holds if and only if $G'$ and $G$ share the same $t$-clique Perron vector.
                By symmetry, in the graph $G'$, the vertices in $M_1'$ and those in $M_0$ have the same $t$-clique Perron vector components,
                however, in the graph $G$, by comparing the components of the vertices in $M_1 $ and $M_0 $ through Equation (1), we have $x_0 > x_1$, which implies that the $t$-clique Perron vectors of $G$ and $G'$ are different, therefore, the equality does not hold. The second inequality holds because $m_2 \geq m_i$ and $x_1 \geq x_i$ for all $i \in {2, 3, \cdots, l}$, hence,

                    \begin{align*}
                        \rho_t(G') > \rho_t(G) = \rho_t(\bar{G}),
                    \end{align*}
                this is a contradiction.
            \end{proof}


            \begin{claim}\label{CC2}
                $H$ have exactly one level-2 subordinate component.
            \end{claim}

            \begin{proof}

                Let $D_1$ and $D_2$ be two different level-2 components which are subordinate components of $H$. Denote the subordinate components of $D_1$, ordered by clique size from largest to smallest, as $M_1^1, M_1^2, \ldots, M_1^{p_1}$, and those of $D_2$ as $M_2^1, M_2^2, \ldots, M_2^{p_2}$.
                Obviously $un(D_i) \leq p_i - 1$; otherwise, by connecting any two vertices in $D_i$ to obtain a  complete graph $D_i'$, the circumference of $G$ does not increase by $p(D_i)=p(D_i')=|V(D_i)|$, and the $t$-clique spectral radius increases, because the operation introduces new $t$-cliques, for instance, for $u \in Un(D_i)$ and $v, w \in V(D_i) \setminus Un(D_i)$ with $v \nsim w$, since $G$ is $t$-clique connected, it follows from the neighborhoods of these vertices that there exists a $t$-clique $K$ in $G$ containing $u$ and $v$, then, connecting $w$ and $v$ yields a new $t$-clique $K'$, where $V(K') = V(K) \setminus \{u\} \cup \{w\}$.

                By Claim~\ref{CC1}, each level-2 component has at most one subordinate component with more than one vertex, therefore, $|M_i^{j_i}| = 1$, where $j_i \in \{2,3,\ldots,p_i\}$ and $i \in \{1,2\}$.

                Let $x$ be the $t$-clique Perron vector of $G$, by symmetry, we may assume that the entry of $x$ corresponding to any vertex in $Un(D_1)$ is $x_1$, and that corresponding to any vertex in $Un(D_2)$ is $x_2$, let the entries corresponding to $M_1^2, \ldots, M_1^{p_1}$ be denoted by $x_1'$, and those corresponding to $M_2^2, \ldots, M_2^{p_2}$ by $x_2'$.

                Now consider modifying the graph $G$ as follows: delete all edges with one incident vertex in $\bigcup_{i=2}^{p_2} V(M_2^i)$ and the other in $Un(D_2)$, and then add all possible edges between vertices in $\bigcup_{i=2}^{p_2} V(M_2^i)$ and $Un(D_1)$ to obtain a new graph $G'$. Similarly, delete all edges with one incident vertex in $\bigcup_{i=2}^{p_1} V(M_1^i)$ and the other in $Un(D_1)$, and then add all possible edges between vertices in $\bigcup_{i=2}^{p_1} V(M_1^i)$ and $Un(D_2)$ to obtain a new graph $G''$, clearly, $c(G') \leq c(G) \leq c(\bar{G})$ and $c(G'') \leq c(G) \leq c(\bar{G})$. It is easy to see that $L_r(G, D_1, x) = L_r(G, D_2, x)$ holds for any $r$, hence
                    \begin{align*}
                        &\rho_t(G')-\rho_t(G)
                        \geq  A_t(G')x^t-A_t(G)x^t
                        \\= & t\sum_{r=1}^{t}L_{t-r}(G, D_1, x) (p_2-1)x_2' \left( \begin{pmatrix}un(D_1)\\r-1\end{pmatrix}x_1^{r-1}-\begin{pmatrix}un(D_2)\\r-1\end{pmatrix}x_2^{r-1} \right),
                    \end{align*}
                and
                    \begin{align*}
                        &\rho_t(G'')-\rho_t(G)
                        \geq  A_t(G'')x^t-A_t(G)x^t
                        \\= & t\sum_{r=1}^{t}L_{t-r}(G, D_1, x) (p_1-1)x_1' \left( \begin{pmatrix}un(D_2)\\r-1\end{pmatrix}x_2^{r-1}-\begin{pmatrix}un(D_1)\\r-1\end{pmatrix}x_1^{r-1} \right),
                    \end{align*}
                By comparing the two equations above, it is clear that either $\rho_t(G')-\rho_t(G)\geq 0$ or $\rho_t(G'')-\rho_t(G)\geq 0$, without loss of generality, assume $\rho_t(G')-\rho_t(G)\geq 0$, where equality holds if and only if the graph $G'$ has the same $t$-clique Perron vector as $G$ and
                    \begin{align*}
                        \begin{pmatrix}un(D_1)\\r-1\end{pmatrix}x_1^{r-1}= \begin{pmatrix}un(D_2)\\r-1\end{pmatrix}x_2^{r-1},
			        \end{align*}
                however, from the characteristic equation corresponding to any vertex in ${Un}(D_1)$ in Equation (1), we conclude that $x$ cannot be an eigenvector of $G'$, thus the equality condition cannot hold. Furthermore, since $G'$ has strictly fewer level-2 subordinate components than $G$, and $\rho_t(G') > \rho_t(G) =\rho_t(\bar{G})$, we arrive at a contradiction.

            \end{proof}

            We now proceed to conclude the proof of the lemma. By Claim~\ref{CC2},
            the subordinate component set of $H$ contains exactly one level-2 component, i.e., $D_1$.
            With the exception of $D_1$, all other subordinate components of $H$ are cliques, ordered by size in descending sequence as $M_1, M_2, \ldots, M_l$, with their respective sizes denoted by $m_1, \ldots, m_l$. Obviously $l \neq 0$, otherwise $H$ would be a level-2 component, which is a contradiction.


            If $m_1 = 1$, then for $i \in [l]$, $M_i$ contains only one vertex, after fully connecting these vertices with the vertices in $Un(D_1)$, the graphs $G'$ and $H'$ are obtained from the graph $G$ and its subgraph $H$, respectively. Then, compared to $G$, $G'$ generates new $t$-cliques and $p(H') = p(H)$, so $\rho_t(G') > \rho_t(G)$ and $c(G') \leq c(G)$.

            If $l = 1$, by adding all edges between the vertex sets $V(Un(D_1))$ and $V(M_1)$, we immediately obtain graph $G'$ with $\rho_t(G') > \rho_t(G)$ and $c(G')\leq c(G)$.

            Therefore, it suffices to consider the case when $m_1 \geq 2$ and $l \geq 2$.

            When $m_1 \geq 2$ and $l \geq 2$, select a maximum clique on $D_1$ and denote it as $K(D_1)$, let $|K(D_1)| = k(D_1)$. It is easy to see that $Un(D_1) \subsetneq K(D_1)$.

            If $un(D_1) > m_1$, similarly to Claim~\ref{CC1}, it can be shown that the following operation will yield a graph $G'$ such that $c(G') \leq c(G)$ and $\rho_t(G') > \rho_t(G)$:

            \noindent (1) Arbitrarily choose $m_1 - 1$ vertices from $V(Un(D_1))$ to form an induced subgraph of $Un(D_1)$, denoted by $Un(D_1)'$. Connect each pair $(i, j)$, where $i \in V(Un(D_1)')$ and $j \in \bigcup_{i=1}^l V(M_i)$;

            \noindent (2) For each $i \in {1, \cdots, l}$, delete all edges in $E(M_i)$.


            If $k(D_1) \geq m_1 \geq un(D_1)$, consider the following method to move the neighbors of the vertices in $M_i$ to the vertices of $K(D_1)$. Denote the vertices in $M_1$ as $v_1^1, v_1^2, \cdots, v_1^{m_1}$.
            Select $m_1 - 1$ vertices on $K(D_1)$ (preferentially choosing vertices from $Un(D_1)$), denoted as $u_1^1, u_1^2, \cdots, u_1^{m_1 - 1}$. Let $H_1^0 = H$. For $i = 1, \cdots, m_1 - 1$, let $H_1^i = {H_1^{i-1}}_{v_1^i \to u_1^i}$, obtaining the graph $H_1^{m_1 - 1}$, then, the induced subgraph of $H_1^{m_1 - 1}$ on $V(Un(H)) \cup V(D_1) \cup V(M_1)$ is a proper subgraph of $S_{|V(H)|, m_1+un(H)-1, k(D_1) - m_1 + 1}$.
            Next, denote the vertices in $M_2$ as $v_2^1, v_2^2, \cdots, v_2^{m_2}$. Let $H_2^0 = H_1^{m_1 - 1}$, for $i = 1, \cdots, m_2 - 1$, let $H_2^i = {H_2^{i-1}}{v_2^i \to u_1^i}$, obtaining the graph $H_2^{m_2 - 1}$.
            Continue this process iteratively for $i = 3, \ldots, l$, applying the same operation to each $M_i$, to finally obtain the graph $H_l^{m_l - 1}$.
            Then, $H_l^{m_l - 1}$ is still a proper subgraph of
                \begin{align*}
                        S_{|V(H)|, m_1+un(H)-1, k(D_1) - m_1 + 1}.
                \end{align*}
            It is easy to see that $p(S_{|V(H)|, m_1+un(H)-1, k(D_1) - m_1 + 1}) = 2un(H)+m_1+k(D_1)-1 \leq p(H)$.
            Furthermore, by replacing the edge subset $E(H)$ in $G$ with
                \begin{align*}
                        E(S_{|V(H)|, m_1+un(H)-1, k(D_1) - m_1 + 1})
                \end{align*}
            we obtain the graph $G'$, we have $c(G') \leq c(G)$. And then, according to Lemma~\ref{CL1}, there is $\rho_t(G') > \rho_t(G)$.


            If $m_1 > k(D_1)$, similar to the case when $k(D_1) \geq m_1 \geq un(D_1)$, we can likewise obtain a graph $G'$ such that $c(G') \leq c(G)$ and $\rho_t(G') > \rho_t(G)$ by sequentially transferring the neighbors of the vertices in $K(D_1)$ and $M_i$ to the vertices in $M_1$ for all $i = {2, 3, \cdots, l}$.



            That is, there exists a $C_{\geq k}$-free graph $G'$ with $n$ vertices such that
                    \begin{align*}
                        \rho_t(G')>\rho_t(G) = \rho_t(\bar{G}),
                    \end{align*}
            This contradicts the assumption that $\bar{G}\in \mathrm{SPEX}_t'(n,C_{\geq k})$.


        \end{proof}

\end{document}
